% !TEX root = ../main.tex
\chapter{Fonaments tècnics}\label{ch:fonaments}
\fancyhead[RE]{\slshape \textsf{Capítol \thechapter.\ Fonaments tècnics}}
\fancyhead[LO]{\slshape \textsf{\nouppercase \rightmark}}

\section{Dades multiespectrals i rang numèric}\label{sec:dades_multi}

Una adquisició multiespectral conté diverses mesures de la mateixa superfície
en intervals de longitud d'ona diferents. A Sentinel-2, per exemple, les bandes
B02, B03, B04 i B08 corresponen aproximadament al blau, verd, vermell i
infraroig proper; altres bandes cobreixen el red edge, vapor d'aigua o SWIR
\cite{esa2015sentinel2}. El producte original no té totes les bandes a la
mateixa resolució espacial. El conjunt de dades utilitzat en aquesta memòria
parteix de retalls ja preparats amb vuit bandes i una graella comuna de
$512\times512$ mostres.

La representació d'entrada és:

\begin{equation}
\mathbf{x}\in\{0,\ldots,65535\}^{B\times H\times W},
\qquad B=8,\;H=W=512.
\label{eq:input_tensor}
\end{equation}

El límit superior $65535$ no és arbitrari: és $2^{16}-1$, el valor màxim d'un
enter sense signe de 16 bits. Els fitxers RAW d'entrada s'emmagatzemen en aquest
format \code{uint16}. Abans d'entrar a la xarxa, les mostres es normalitzen a
l'interval aproximat $[0,1]$ dividint per aquest mateix valor, tal com es va
corregir i validar en la implementació C del treball previ
\cite{anez2026treball}.

La mida crua d'un retall és:

\begin{equation}
S_{\mathrm{raw}}=BHW\cdot2
=8\cdot512\cdot512\cdot2
=4\,194\,304\;\mathrm{bytes}.
\label{eq:raw_size}
\end{equation}

En aquesta memòria, quan s'utilitza \emph{bps}, es mesuren bits per mostra
espectral, és a dir, el denominador és $BHW$. Això és important perquè una
imatge multiespectral no conté només píxels espacials: conté píxels per banda.

La compressió és possible perquè hi ha redundància. La correlació espacial
permet predir una mostra a partir de les veïnes. La correlació espectral permet
expressar una banda a partir de combinacions d'altres bandes. SORTENY tracta
aquestes dues fonts de redundància amb transformades apreses.

\section{SORTENY des del punt de vista de la implementació}
\label{sec:arquitectura}

Després de situar SORTENY en el context dels còdecs neuronals per a observació
de la Terra, cal concretar quins elements interns fan possible el control de
qualitat que s'estudia en aquesta memòria. En particular, aquest capítol
descriu la geometria de la xarxa, les operacions amb més cost computacional, el
paper del modulador de qualitat i les mètriques que s'utilitzaran per avaluar
els resultats.

El pipeline conceptual és:

\begin{equation}
\mathbf{x}\xrightarrow{g_a}
\mathbf{y}\xrightarrow{\mathcal{Q}_{\lambda}}
\hat{\mathbf{y}}\xrightarrow{g_s}\hat{\mathbf{x}}.
\label{eq:sorteny_pipeline}
\end{equation}

$g_a$ és la transformada d'anàlisi, $\mathbf{y}$ són els latents,
$\mathcal{Q}_{\lambda}$ és la quantització condicionada pel paràmetre de
qualitat i $g_s$ és la transformada de síntesi. Els pesos ja estan entrenats.
El treball d'aquesta memòria es centra en executar i controlar aquesta
inferència, no en reentrenar la xarxa.

\subsection{Transformada espectral}

La transformada espectral combina les $B$ bandes en cada posició espacial. En
forma simplificada:

\begin{equation}
\mathbf{z}_{h,w}=\mathbf{W}_{s}\mathbf{x}_{h,w}.
\label{eq:spectral_transform}
\end{equation}

$\mathbf{W}_{s}$ és una matriu apresa. A diferència d'una transformada fixa,
pot adaptar-se a l'estadística del conjunt d'entrenament. Computacionalment,
l'operació és independent entre posicions espacials: cada píxel pot processar-se
separadament, mantenint l'ordre intern de suma sobre bandes. Aquesta propietat
és rellevant quan es paral·lelitza el C sense alterar el resultat numèric.

\subsection{Convolucions i cost principal}

Les convolucions redueixen progressivament la resolució espacial i concentren
informació en canals latents. Una convolució amb kernel $K\times K$ pot
escriure's com:

\begin{equation}
y_{c_o,h,w}=b_{c_o}+
\sum_{c_i}\sum_{u=0}^{K-1}\sum_{v=0}^{K-1}
W_{c_o,c_i,u,v}\,x_{c_i,h+u,w+v}.
\label{eq:conv2d}
\end{equation}

La fórmula mostra per què és una operació costosa. Per cada mostra de sortida
cal recórrer canals d'entrada i posicions del kernel. En SORTENY, els kernels
espacials són de $5\times5$ i el nombre de canals intermedis és elevat. Per
això, el temps total del còdec està dominat per aquestes operacions i no pel
càlcul de regles semàntiques senzilles com un índex espectral.

\subsection{GDN i IGDN}

Entre convolucions s'utilitza \emph{Generalized Divisive Normalization} (GDN),
una normalització no lineal habitual en compressió neuronal
\cite{balle2017end}. Per una posició espacial fixa:

\begin{equation}
y_i=\frac{x_i}
{\sqrt{\beta_i+\sum_j\gamma_{ij}x_j^2}}.
\label{eq:gdn}
\end{equation}

El canal $i$ no depèn només d'ell mateix, sinó també de la suma ponderada dels
altres canals $j$. La IGDN aplica la transformació complementària al decoder.
A diferència d'una activació simple com ReLU, la GDN no només talla o escala
valors de manera independent: normalitza cada resposta segons l'energia dels
altres canals. En compressió neuronal això és útil perquè tendeix a reduir
dependències estadístiques entre coeficients latents, un efecte sovint descrit
com una aproximació a ``gaussianitzar'' o fer més regular la distribució dels
coeficients. Uns latents menys dependents són més adequats per al model de
probabilitat i per al codificador per entropia posterior.

Aquest guany de modelatge té un cost computacional: cal calcular sumes entre
canals per moltes posicions espacials. Aquesta és una de les raons per les quals
l'optimització C ha de respectar l'ordre de les operacions i validar paritat
numèrica.

\subsection{Modulador de qualitat}

SORTENY és multiràtio perquè incorpora una subxarxa de modulació condicionada
pel paràmetre $\lambda$ \cite{mijares2023reduced,mijares2025sorteny}. El
modulador genera factors que escalen els latents abans de quantitzar-los i que
el decoder necessita per invertir l'operació de manera coherent.

És important no confondre $M(\lambda)$ amb una màscara ROI. $M(\lambda)$ és un
tensor de factors de modulació après, amb la mateixa funció conceptual per a
tota la imatge o per cada posició latent. Una màscara ROI, en canvi, és una
decisió externa sobre quines zones es volen prioritzar. El Q-map que es defineix
més endavant tradueix aquesta decisió externa a valors de $\lambda$ quantitzats,
però el modulador continua sent el mecanisme intern de la xarxa.

\section{Geometria interna i granularitat del Q-map}
\label{sec:tensor_geometry}

La geometria dels tensors explica dues propietats centrals: el cost de memòria
i la mida de la graella de control de qualitat. La transformada espectral opera
a resolució completa. Després, cada component passa per quatre convolucions de
stride 2. Cada stride 2 redueix l'alçada i l'amplada a la meitat:

\begin{figure}[H]
\centering
\includegraphics[width=0.92\textwidth]{ch3/tensor_geometry_funnel.pdf}
\caption{Reducció espacial de SORTENY fins a la graella latent. La figura
mostra dimensions conceptuals; la implementació reutilitza buffers i no conserva
tots els tensors simultàniament.}
\label{fig:tensor_geometry_funnel}
\end{figure}

\begin{table}[H]
\centering
\small
\begin{tabular}{lrrrrl}
\toprule
Etapa & Canals & Alçada & Amplada & Valors & Operació posterior \\
\midrule
Entrada per banda & 1 & 512 & 512 & 262.144 & Conv. 0 \\
Anàlisi 0 & 128 & 256 & 256 & 8.388.608 & GDN \\
Anàlisi 1 & 128 & 128 & 128 & 2.097.152 & GDN \\
Anàlisi 2 & 128 & 64 & 64 & 524.288 & GDN \\
Latent per component & 384 & 32 & 32 & 393.216 & Modulació i Q \\
Latent complet & $8\cdot384$ & 32 & 32 & 3.145.728 & Codificació \\
\bottomrule
\end{tabular}
\caption{Dimensions principals derivades dels pesos i de la implementació C.
La columna \emph{Valors} indica el nombre d'elements del tensor, no la memòria
total reservada durant tota la inferència.}
\label{tab:tensor_geometry}
\end{table}

El màxim tensor intermedi d'una banda conté més de vuit milions de
\code{float32}, uns 32~MiB. Conservar còpies independents per capa o per fil
superaria ràpidament la memòria disponible en una plataforma de baix recurs.
Aquesta geometria justifica la reutilització de buffers que s'explica al
Capítol~\ref{ch:pipeline}.

\subsection{\texorpdfstring{De $512\times512$ a $32\times32$}{De 512x512 a 32x32}}

Quatre reduccions per factor 2 donen un factor total $2^4=16$. Per tant:

\begin{equation}
\frac{512}{16}=32.
\label{eq:qmap_grid_size}
\end{equation}

El control espacial de qualitat opera sobre aquesta graella latent de
$32\times32$ posicions. De manera pràctica, es parla d'un valor Q per bloc
$16\times16$ de la imatge d'entrada. Aquesta frase descriu la graella de
control, però no vol dir que cada bloc sigui independent visualment.

\begin{figure}[H]
\centering
\includegraphics[width=0.82\textwidth]{ch3/qmap_grid.pdf}
\caption{Correspondència entre la imatge d'entrada i el Q-map. Cada valor Q
controla una cel·la latent associada a un pas de 16 píxels a l'entrada.}
\label{fig:qmap_grid}
\end{figure}

\subsection{Camp receptiu i efecte sobre ROIs petites}

Una cel·la latent no depèn només dels $16\times16$ píxels que li corresponen
per pas espacial. Les convolucions amb kernel $5\times5$ fan que cada sortida
vegi un veïnat de l'entrada. Si $r_\ell$ és la mida del camp receptiu després
de la capa $\ell$ i $j_\ell$ és el salt equivalent entre posicions d'aquella
capa respecte l'entrada:

\begin{equation}
r_0=1,\qquad j_0=1,
\end{equation}

\begin{equation}
r_{\ell+1}=r_\ell+(5-1)j_\ell,\qquad
j_{\ell+1}=2j_\ell.
\label{eq:receptive_field}
\end{equation}

Això dona:

\begin{table}[H]
\centering
\small
\begin{tabular}{lrrrrr}
\toprule
Després de la capa & 0 & 1 & 2 & 3 & 4 \\
\midrule
Salt equivalent $j$ & 1 & 2 & 4 & 8 & 16 \\
Camp receptiu $r$ & 1 & 5 & 13 & 29 & 61 \\
\bottomrule
\end{tabular}
\caption{Creixement del camp receptiu amb quatre convolucions $5\times5$ i
stride 2.}
\label{tab:receptive_field}
\end{table}

La conseqüència és important per a la qualitat espacial: una ROI petita o molt
dispersa pot quedar afectada pels blocs veïns. Encara que la cel·la latent
marcada com a ROI rebi una qualitat alta, part de la informació que contribueix
a reconstruir-la pot haver passat per cel·les properes amb una qualitat diferent.

\begin{figure}[H]
\centering
\includegraphics[width=0.78\textwidth]{ch3/receptive_field_overlap.pdf}
\caption{Diferència entre graella de control i camp d'influència. El bloc de
control és de 16 píxels, però el camp receptiu teòric arriba aproximadament a
61 píxels.}
\label{fig:receptive_field_overlap}
\end{figure}

\section{Taxa, distorsió i qualitat}\label{sec:metriques_teoria}

\subsection{MSE, MAE i PSNR}

Per a $N=BHW$ mostres, l'error quadràtic mitjà és:

\begin{equation}
\mathrm{MSE}(\mathbf{x},\hat{\mathbf{x}})
=\frac{1}{N}\sum_{i=1}^{N}(x_i-\hat{x}_i)^2.
\label{eq:mse}
\end{equation}

L'error absolut mitjà és:

\begin{equation}
\mathrm{MAE}(\mathbf{x},\hat{\mathbf{x}})
=\frac{1}{N}\sum_{i=1}^{N}|x_i-\hat{x}_i|.
\label{eq:mae}
\end{equation}

Per dades de 16 bits, el PSNR es calcula amb el mateix màxim dinàmic que a
l'Equació~\ref{eq:input_tensor}:

\begin{equation}
\mathrm{PSNR}=10\log_{10}
\left(\frac{65535^2}{\mathrm{MSE}}\right)\;\mathrm{dB}.
\label{eq:psnr}
\end{equation}

Un PSNR més alt implica menys MSE. La diferència en dB és logarítmica; no es
pot interpretar linealment com un percentatge de píxels correctes. El PSNR és
especialment adequat en aquest treball perquè el model fixed-quality que
s'utilitza com a base científica està formulat en termes de MSE
\cite{mijares2024fixed}. Altres mètriques, com SSIM, modelen estructura visual
\cite{wang2004ssim}, però no són la variable que calibra el mètode.

\subsection{Bitrate i fitxer \code{.tfci}}

El bitrate es mesura en bits per mostra espectral:

\begin{equation}
R=\frac{8S}{BHW}\;\bps,
\label{eq:bps_general}
\end{equation}

on $S$ és la mida codificada en bytes. Quan el fitxer és el resultat del
pipeline de TensorFlow Compression, s'utilitza l'extensió \code{.tfci}. Aquest
fitxer conté latents codificats i la informació lateral que el decoder de la
referència necessita per reconstruir la imatge. Per això és la mesura que
s'utilitza com a bitrate real en les validacions CSMR.

\begin{equation}
R_{\mathrm{tfci}}=
\frac{8S_{\mathrm{tfci}}}{BHW}\;\bps.
\label{eq:bps}
\end{equation}

El bitstream C d'aquesta memòria té una funció diferent: empaqueta capçalera,
Q-map i latents per validar la ruta C, però encara no incorpora el mateix
codificador per entropia o codificador aritmètic que la referència. Per aquest
motiu, la mida C directa no s'interpreta com a bitrate final de compressió.

Si un mètode A té bitrate $R_A$ i el baseline B té $R_B$, l'estalvi relatiu és:

\begin{equation}
\mathrm{estalvi}_{A/B}
=100\left(1-\frac{R_A}{R_B}\right)\%.
\label{eq:saving}
\end{equation}

Un estalvi del 10\% significa que A envia el 90\% dels bits de B. No equival a
dir que la compressió s'ha duplicat.

\subsection{Entropia i proxies}

L'entropia empírica d'un alfabet latent és:

\begin{equation}
H=-\sum_k p_k\log_2 p_k.
\label{eq:entropy}
\end{equation}

Una entropia menor suggereix que els símbols podrien codificar-se amb menys
bits, però no especifica el resultat d'un codificador concret
\cite{cover2006information}. També es pot aplicar una compressió generalista
com zlib per obtenir un indici ràpid de compresibilitat. En aquesta memòria,
entropia i zlib són proxies de selecció: ajuden a decidir quins casos mereixen
una validació completa, però no substitueixen el bitrate \code{.tfci}.

\section{Qualitat espacial, ROI i fons}\label{sec:roi_metrics}

Una regió d'interès, o ROI, és el conjunt de blocs que es vol protegir amb més
qualitat. En aquest capítol la definició és deliberadament genèrica: la màscara
pot venir d'un operador manual, d'un criteri espectral o d'un altre mecanisme.
Els capítols següents definiran els mecanismes concrets utilitzats.

Sigui $m_{u,v}\in\{0,1\}$ la màscara sobre la graella latent
$32\times32$. Els blocs amb $m=1$ formen la ROI i els blocs amb $m=0$ formen el
fons. Per mesurar PSNR sobre la imatge d'entrada, aquesta màscara de blocs
s'expandeix a les mostres espacials corresponents.

Per una reconstrucció A i un baseline B, sempre amb la mateixa màscara:

\begin{align}
\Delta\mathrm{PSNR}_{\mathrm{ROI}}
&=\mathrm{PSNR}_{\mathrm{ROI}}(A)
-\mathrm{PSNR}_{\mathrm{ROI}}(B),\\
\Delta\mathrm{PSNR}_{\mathrm{fons}}
&=\mathrm{PSNR}_{\mathrm{fons}}(A)
-\mathrm{PSNR}_{\mathrm{fons}}(B).
\label{eq:roi_bg_deltas}
\end{align}

Usar la mateixa màscara és essencial. Si el baseline i el mètode avaluat es
mesuressin sobre regions diferents, els deltes no compararien la mateixa
població de mostres. Una política espacial útil ha de conservar o millorar la
ROI respecte del baseline escollit, mentre pot degradar el fons de manera
controlada si això redueix el bitrate.

La cobertura de la ROI és:

\begin{equation}
C_{\mathrm{ROI}}=
\frac{\sum_{u,v}m_{u,v}}{32\cdot32}\cdot100\%.
\label{eq:roi_coverage}
\end{equation}

Els extrems són poc informatius. Si $C_{\mathrm{ROI}}=0$, no hi ha cap regió
protegida. Si $C_{\mathrm{ROI}}=100\%$, tota la imatge és ROI i no queda fons
d'on obtenir estalvi. Entre aquests extrems, la mida i la forma de la ROI
condicionen la utilitat del control espacial.

\section{Model fixed-quality}\label{sec:model_fq}

El mètode fixed-quality de Mijares i Verdú \etal\ parteix d'una pregunta
inversa: en lloc de provar diferents qualitats fins trobar-ne una adequada,
estima quin $\lambda$ cal utilitzar per assolir un error objectiu
\cite{mijares2024fixed}.

\subsection{De PSNR objectiu a MSE objectiu}

Si l'usuari defineix una qualitat objectiu en PSNR, primer cal convertir-la a
MSE. Reordenant l'Equació~\ref{eq:psnr}:

\begin{equation}
\widehat{\mathrm{MSE}}=
\frac{65535^2}
{10^{\mathrm{PSNR}_{\mathrm{obj}}/10}}.
\label{eq:target_mse}
\end{equation}

El símbol $\widehat{\mathrm{MSE}}$ indica l'error que es voldria obtenir. Aquest
valor no determina directament el byte d'un Q-map. Abans cal tenir en compte
l'error base de la imatge o del bloc.

\subsection{MSE base i modulador mitjà}

El paper defineix $\mathrm{MSE}_0(\mathbf{x})$ com l'error base o mínim associat
a la imatge $\mathbf{x}$ en el model considerat. Aquest valor depèn del
contingut: dues imatges comprimides amb el mateix $\lambda$ poden donar PSNR
diferents perquè no tenen la mateixa dificultat de reconstrucció.

La magnitud mitjana del modulador es modela com una funció lineal:

\begin{equation}
\bar{M}(\lambda)=a\lambda+b.
\label{eq:mod_linear}
\end{equation}

El paper observa que l'error de reconstrucció creix de manera aproximadament
quadràtica amb l'invers d'aquest modulador mitjà. La relació s'escriu com:

\begin{equation}
\widehat{\mathrm{MSE}}
=\mathrm{MSE}_0+
\alpha\,\mathrm{MSE}_0
\left(\frac{1}{2\bar{M}(\lambda)}\right)^2.
\label{eq:mse_model}
\end{equation}

Dividint per $\mathrm{MSE}_0$, es veu millor la idea de l'error relatiu:

\begin{equation}
\frac{\widehat{\mathrm{MSE}}}{\mathrm{MSE}_0}
=1+
\alpha
\left(\frac{1}{2\bar{M}(\lambda)}\right)^2.
\label{eq:mse_relative_model}
\end{equation}

Aquesta expressió és una aproximació estadística calibrada. No diu que tots els
píxels tinguin el mateix error, sinó que permet estimar quin $\lambda$ produeix
un error mitjà esperat.

\subsection{Estimació de \texorpdfstring{$\lambda$}{lambda}}

Combinant les Equacions~\ref{eq:mod_linear} i~\ref{eq:mse_model}, el paper
despeja el valor de $\lambda$ necessari:

\begin{equation}
\hat{\lambda}
=\frac{1}{a}
\left(
\sqrt{
\frac{\alpha\,\mathrm{MSE}_0(\mathbf{x})}
{4\left(\widehat{\mathrm{MSE}}-\mathrm{MSE}_0(\mathbf{x})\right)}
}
-b
\right).
\label{eq:lambda_from_mse}
\end{equation}

La cadena completa és, per tant:

\begin{figure}[H]
\centering
\includegraphics[width=0.92\textwidth]{ch3/fixed_quality_chain.pdf}
\caption{Cadena conceptual del model fixed-quality: d'una qualitat objectiu es
dedueix un MSE objectiu, després un $\lambda$ estimat i finalment una
representació quantitzada per transmetre al decoder.}
\label{fig:fixed_quality_chain}
\end{figure}

L'Equació~\ref{eq:lambda_from_mse} només té sentit si el target és assolible:
el denominador ha de ser positiu i el resultat ha de quedar dins el rang de
$\lambda$ per al qual el model i la calibració són vàlids. Quan no passa, el
sistema ha de saturar al límit disponible o marcar el cas com no assolible.

\subsection{Quantització de \texorpdfstring{$\lambda$}{lambda} a Q}

Un array de $\lambda$ en coma flotant seria informació lateral costosa. El paper
quantitza cada $\lambda$ a un enter de 8 bits:

\begin{equation}
Q=
\operatorname{round}\left[
255\frac{\lambda-\lambda_{\min}}
{\lambda_{\max}-\lambda_{\min}}
\right],
\qquad Q\in\{0,\ldots,255\}.
\label{eq:lambda_to_q}
\end{equation}

El decoder pot recuperar el valor aproximat amb:

\begin{equation}
\hat{\lambda}=
\lambda_{\min}+
\frac{Q}{255}
(\lambda_{\max}-\lambda_{\min}).
\label{eq:q_to_lambda}
\end{equation}

Per tant, el byte Q no és una quantització clàssica directa dels píxels. És una
quantització lineal del paràmetre $\lambda$, que després actua a través del
modulador de la xarxa. En aquesta arquitectura, valors Q més alts corresponen a
valors $\lambda$ més alts i, en general, a més qualitat i més taxa.

\subsection{Cost de la informació lateral}

Quan la qualitat varia espacialment, el decoder ha de conèixer quina qualitat
s'ha aplicat a cada posició latent. Aquesta informació no forma part de la
imatge reconstruïda, però és necessària per poder descodificar-la correctament;
per això s'anomena \emph{informació lateral}. En la ruta C d'aquesta memòria,
aquesta informació lateral és el Q-map escrit explícitament al bitstream
intermedi.

Amb una graella $32\times32$, un Q-map conté 1.024 valors. Si cada valor és un
byte, la informació lateral ocupa 1.024 bytes. Expressada en bits per píxel
espacial, la mateixa convenció que utilitza el paper fixed-quality per comparar
arrays de $\lambda$, aquesta mida és:

\begin{equation}
R_{Q,\mathrm{espacial}}=
\frac{8\cdot1024}{512\cdot512}
=0.03125\;\text{bits/píxel}.
\label{eq:qmap_side_spatial_bps}
\end{equation}

La resta de resultats d'aquesta memòria, però, es donen en bits per mostra
espectral. Com que cada retall té vuit bandes, el mateix cost es reparteix sobre
$8\cdot512\cdot512$ mostres:

\begin{equation}
R_Q=
\frac{8\cdot1024}{8\cdot512\cdot512}
=0.00390625\;\bps.
\label{eq:qmap_side_bps}
\end{equation}

La diferència entre $0.03125$ i $0.00390625$ no és una contradicció; només
canvia el denominador. La informació lateral és petita respecte del volum RAW,
però és real i s'ha de transportar o reconstruir de manera equivalent al
decoder.

\section{Lectura dels capítols següents}\label{sec:criteris_lectura}

Amb aquests fonaments es pot separar el problema en tres decisions. Primer,
quin valor de qualitat es vol assignar. Segon, si aquest valor és global o pot
variar sobre la graella $32\times32$. Tercer, quin criteri decideix quins blocs
mereixen més o menys qualitat.

Els capítols següents concreten aquestes decisions en una implementació C i en
experiments. El Capítol~\ref{ch:pipeline} descriu el còdec i el bitstream. El
Capítol~\ref{ch:control} defineix els modes de Q-map i els presets automàtics.
El Capítol~\ref{ch:metodologia} explica com es mesuren resultats amb màscares
comunes, C-only i CSMR. Finalment, els capítols de resultats indiquen quins
compromisos taxa--distorsió s'obtenen i quin cost tenen en una plataforma de
baix recurs.
